% ---------------------------------------------------------------------------
% Fachmodul Disjunktive Normalform (Duy Tan Pham)
%
% Inhaltlich abgeleitet aus Abschnitt 4.6 der Dokumentation
% (includes/04_Designentscheidungen/04_06.tex).
% ---------------------------------------------------------------------------

\addsectionframe{Fachmodul: Disjunktive Normalform}

% ---------------------------------------------------------------------------
\begin{itemframe}{Disjunktive Normalform -- Überblick}
  \item Zwei Aufgabenmodi: Ausdruck $\rightarrow$ Wahrheitstabelle und Wahrheitstabelle $\rightarrow$ DNF
  \item Drei Ergebnisstufen: kanonische DNF, minimale DNF und optional global kleinste UND/ODER-Formel
  \item Minimierung mit Quine-McCluskey und Branch-and-Bound
  \item Semantische Prüfung erkennt logische Äquivalenz
  \item Begrenzung auf maximal 4 Variablen für exakte Synthese
  \note{
    \textbf{Person: Tan}

    \textbf{Zeit:  1 min}

    \begin{itemize}
      \item \textbf{Aufbau:} UND-Terme werden durch ODER verbunden
      \item \textbf{Zwei Wege:} Ausdruck zur Wahrheitstabelle oder Wahrheitstabelle zur DNF
      \item \textbf{Kanonisch:} Jede 1-Zeile ergibt einen vollständigen Term mit allen Variablen
      \item \textbf{Minimal:} Weniger Literale, exakt mit QMC + B\&B
      \item \textbf{Global:} Ausklammern erlaubt, daher nicht zwingend DNF
      \item \textbf{Äquivalenz:} Funktion über Wahrheitstabelle vergleichen
      \item \textbf{Grenze:} Bis vier Variablen garantiert optimal
    \end{itemize}
  }
\end{itemframe}

% ---------------------------------------------------------------------------
\begin{contentsframe}{Schrittausführung}
\centering
\vspace{1mm}
\resizebox{0.95\textwidth}{!}{%
\begin{tikzpicture}[node distance=15mm and 11mm]
  \node[diagramnode, minimum width=3.3cm, font=\normalsize] (input) {Ausgangsdaten:\\Formel oder Tabelle};
  \node[diagramnode, right=of input, minimum width=3.3cm, font=\normalsize] (ast) {Syntaxprüfung\\Parser \& AST};
  \node[diagramnode, right=of ast, minimum width=3.5cm, font=\normalsize] (table) {Referenztabelle\\1-Zeilen bestimmen};
  \node[diagramnode, right=16mm of table, minimum width=3.5cm, font=\normalsize] (cdnf) {Kanonische DNF\\(Vollminterme)};

  \node[diagramnode, below=15mm of input, minimum width=3.3cm, font=\normalsize] (student) {Antwort\\je nach Modus};
  \node[diagramnode, below=15mm of ast, minimum width=3.3cm, font=\normalsize] (check) {Prüfung: Werte, Form\\und Äquivalenz};
  \node[diagramnode, below=15mm of table, minimum width=3.8cm, font=\normalsize] (exact) {Global kleinste Form\\exakt bis 4 Variablen};
  \node[diagramnode, below=15mm of cdnf, minimum width=3.5cm, font=\normalsize] (qmc) {Minimale DNF\\QMC + B\&B};

  \draw[diagramarrow] (input) -- node[above, font=\small, text=black] {Formel} (ast);
  \draw[diagramarrow] (ast) -- (table);
  \draw[diagramarrow] (input.north) to[out=70, in=110]
    node[above, font=\small, text=black] {Tabelle} (table.north);
  \draw[diagramarrow] (table) -- node[above, font=\small, text=black] {1-Zeilen} (cdnf);
  \draw[diagramarrow] (cdnf) -- (qmc);
  \draw[diagramarrow] (table) --
    node[pos=0.52, right, xshift=2mm, font=\small, text=black] {Funktion} (exact);
  \draw[diagramarrow, dashed] (qmc) -- node[above, font=\small, text=black] {Obergrenze} (exact);

  \draw[diagramarrow] (student) -- (check);
  \draw[diagramarrow] (table) --
    node[pos=0.68, above, xshift=-2mm, font=\small, text=black] {Referenz} (check);
  \draw[diagramarrow] (check.south) to[out=215, in=325] node[below, font=\small, text=black] {Feedback \& Äquivalenz} (student.south);
\end{tikzpicture}%
}
\note{
  \textbf{Person: Tan}

  \textbf{Zeit:  1,5 min}

  \begin{itemize}
    \item \textbf{Ausdrucksmodus:} Referenztabelle erzeugen und eingegebene Werte vergleichen
    \item \textbf{Tabellenmodus:} Vorgegebene Tabelle direkt als Referenz verwenden
    \item \textbf{Gemeinsamer Weg:} 1-Zeilen werden zur kanonischen DNF
    \item \textbf{Minimal:} Primimplikanten bilden und exakt auswählen
    \item \textbf{Global:} Minimale DNF als Obergrenze, kleinere Formel exakt suchen
    \item \textbf{Feedback:} Form und logische Äquivalenz mit der Referenztabelle prüfen
  \end{itemize}
}
\end{contentsframe}
